\section{Einleitung}

Die Grundlagen der Kinematik stellen einen wesentlichen Bestandteil der Robotik dar und sind sowohl für mobile als auch für stationäre Roboter von zentraler Bedeutung.
Während es bei stationären Robotern häufig darum geht, die Position des Endeffektors präzise zu berechnen,
konzentriert sich die Kinematik bei mobilen Robotern vor allem darauf, deren Position zuverlässig bestimmen zu können.
Ein fundiertes Verständnis der erforderlichen kinematischen Modelle ist unerlässlich, um diese Größen korrekt zu berechnen.
\cite[57]{autonomous_mobile_robots}

Im Rahmen dieser Projektarbeit wird ein Framework entwickelt, das die Visualisierung kinematischer Kennwerte ermöglicht.
Das Ziel dieses Frameworks ist es, unterstützend in der Lehre eingesetzt zu werden, um das Verständnis der Kinematik zu fördern und die praktische Anwendung dieser Konzepte zu erleichtern.
Ein besonderer Fokus dieser Arbeit liegt auf der Rotation von Körpern im Raum. Dabei wird insbesondere erläutert,
was unter Eulerwinkeln sowie Roll-Pitch-Yaw-Winkeln zu verstehen ist und worin die Unterschiede zwischen beiden Darstellungen bestehen.
Darüber hinaus werden auch Rotationsmatrizen sowie Quaternionen behandelt, um die verschiedenen Ansätze zur Beschreibung von Rotationen umfassend darzustellen.